\section{Introduction}

A survey programme covering many countries, regions or domains produces one
estimate of some summary per population per round. A second analysis then uses
those estimates to say something about a population the survey did not cover:
a region added to a sampling frame, a country joining a programme, a domain too
small to have been sampled. Call the target $\theta_{K+1}$, the latent summary of
that population, and the observed inputs $Y_1,\dots,Y_K$, the survey estimates
for the populations that were covered.

The obstacle is that each $Y_c$ carries sampling error. The spread of the $Y_c$
is larger than the spread of the $\theta_c$, so an interval calibrated on the
observed spread is too wide for the latent target. Removing the difference means
subtracting a design variance from a between-population variance, and the design
variances are themselves estimated---from stratified cluster designs with, in the
European Social Survey regions we analyse, a median of 22 to 29 degrees of
freedom per population.

Two literatures meet here and neither answers the question. Small-area
estimation treats the sampling variances as known and targets areas that were
sampled \citep{fay1979estimates,rao2015small}; where it treats them as estimated
it corrects the mean squared error of an \emph{in-sample} predictor
\citep{wang2003mse,rivest2003mse,you2006hierarchical,maiti2014,sugasawa2017}.
Recent conformal work reaches valid confidence intervals for in-sample areas
without estimating the sampling variances at all: \citet{zhang2026conformal}
target ``the unknown expectations of the in-sample outcomes \dots{} conditional
on the realised sample'', and their split conformal construction is, in their
words, ``easier to implement than the direct or EBLUP interval, because there is
no need at all to estimate the sampling variances'' (their \S3.4). That is an
argument for their target, not against ours: an interval for a population the
survey never covered cannot avoid separating the two scales. Conformal
prediction of a future \emph{unit} within an area \citep{bersson2024optimal} is
a third target again. Prediction for an out-of-sample area is standard
through synthetic estimation, and the linear mixed-model literature does supply
prediction intervals for targets containing a new random effect
\citep{chatterjee2008}---but with the sampling variances treated as known.

\paragraph{What this paper adds.} Not another way to smooth design variances, and
not a claim that conformal prediction is robust here; our exact inference rests
on a normal--$\chi^2$ model and we say so. The contribution is in two parts.

\emph{A result about the problem.} Section~\ref{sec:sensitivity} isolates the
scale elasticity and shows that it, rather than the quality of the scale
estimate, decides which construction is narrowest. The consequences are checkable
and we check them: the number of populations required is smaller by $\rho^{-4}$
than a criterion on the variance component implies; two constructions cross, at a
computable point; and for one matched pair the scale bound cancels entirely,
giving a width ratio that depends only on $K$ and the error budget. That last
prediction was registered before the experiment that tests it, and one of its
three registered claims was refuted---the constant, by an amount we locate.

\emph{A procedure that exploits it.} Section~\ref{sec:method} constructs an
interval whose sensitivity to the unknown scale is the damped one, carries the
uncertainty of an estimated variance structure rather than substituting a fitted
value, and is computable by a conservative bound rather than a search.

\paragraph{What we do not claim.} The guarantee is proved in a restricted normal
model. Under a complex design its conditions fail, measurably, and we report
where the procedure nonetheless behaves and where it does not. On real survey
data we report widths, inputs and diagnostics, and no coverage: the latent target
is unobserved and a held-out population's direct estimate is not a substitute for
it.
